\documentclass{article}

\usepackage{amsmath}
\usepackage{amssymb}
\usepackage{xcolor}
\usepackage{tikz}

\usepackage[utf8]{inputenc}
\usepackage[T1]{fontenc}

\usepackage[polish]{babel}

\usepackage{listings}
\lstset{
    language=SQL,
    inputencoding=utf8x,
    extendedchars=\true,
    literate={ą}{{\k{a}}}1
             {Ą}{{\k{A}}}1
             {ę}{{\k{e}}}1
             {Ę}{{\k{E}}}1
             {ó}{{\'o}}1
             {Ó}{{\'O}}1
             {ś}{{\'s}}1
             {Ś}{{\'S}}1
             {ł}{{\l{}}}1
             {Ł}{{\L{}}}1
             {ż}{{\.z}}1
             {Ż}{{\.Z}}1
             {ź}{{\'z}}1
             {Ź}{{\'Z}}1
             {ć}{{\'c}}1
             {Ć}{{\'C}}1
             {ń}{{\'n}}1
             {Ń}{{\'N}}1
}

\title{Grupy i Ciała}
\date{\today}

\begin{document}
\maketitle

%%%%%%%%%%%%%%%%%%%%%%%%%%%%%%%%%%%%%%%%%%%%%%%%%%%%%%%%%%%%%%%
\section{Działanie wewnętrzne}
Dwuargumentowym działaniem wewnętrznym w zbiorze $X$ nazywamy funkcję $h : X \times X \rightarrow X$.
Jeśli $x, y \in X$, to $h(x, y) \in X$ nazywamy wynikiem działania $h$ na argumentach $x$ i $y$.
Wygodniej jest zapisywać $x * y$ zamiast $h(x, y)$.

\subsection{Przykłady}
\begin{itemize}
    \item Dodawanie ($+$) i mnożenie ($\cdot$) są działaniami wewnętrznymi w zbiorach $\mathbb{N}, \mathbb{Z}, \mathbb{W}, \mathbb{R}$.
    \item Działanie $m \bigoplus n = m^n$ jest działaniem wewnętrznym w $\mathbb{N}$.
    \item Odejmowanie nie jest działaniem wewnętrznym w zbiorze $\mathbb{N}$.
    \item Zbiór obrotów kwadratu wokół jego środka przekształcających go na siebie.
\end{itemize}

\subsection{Wybrane właściwości}
Niech $*$ będzie działaniem wewnętrznym na zbiorze $X$.
\begin{description}
    \item[Działanie przemienne] - działanie $*$ nazywamy działaniem przemiennym, gdy zachodzi $\forall_{x,y \in X} x * y = y * x$.
    \item[Działanie łączne] - działanie $*$ nazywamy działaniem łącznym, gdy zachodzi $\forall_{x,y,z \in X} x * (y * z) = (x * y) * z$.
    \item[Element neutralny] - element $e \in X$ nazywamy elementem neutralnym działania $*$ jeśli $\forall_{x \in X} x * e = e * x = x$. Jeżeli działanie posiada element neutralny, to jest on jedyny.
    \item[Element przeciwny (odwrotny, symetryczny)] - element $\tilde{x} \in X$ nazywamy elementem przeciwnym dla elementu $x \in X$ jeśli $x * \tilde{x} = \tilde{x} * x = e$.
Jeżeli działanie posiada element neutralny, to element przeciwny do danego elementu jest wyznaczony jednoznacznie.
\end{description}

\subsection{Działanie addytywne i multiplikatywne}
\begin{itemize}
    \item Jeżeli działanie $*$ ma właściwości podobne do dodawania, to nazywamy je addytywnym i zwykle używamy symbolu $+$ lub $\bigoplus$.
    Jego element neutralny $e$ nazywamy "zerem" i oznaczamy $\textbf{0}$, natomiast element przeciwny do $x$ oznaczamy $-x$.
    \item Jeżeli działanie $*$ ma właściwości podobne do mnożenia, to działanie nazywamy multiplikatywnym i oznaczamy zwykle $\cdot$ lub $\bigotimes$.
    Element neutralny $e$ takiego działania nazywamy "jedynką" i oznaczamy $\textbf{1}$, natomiast element odwrotny $\tilde{x}$ do elementu $x$ oznaczamy $x^{-1}$.
\end{itemize}

%%%%%%%%%%%%%%%%%%%%%%%%%%%%%%%%%%%%%%%%%%%%%%%%%%%%%%%%%%%%%%%
\section{Działanie zewnętrzne}
Dwuargumentowym działaniem zewnętrznym w zbiorze $X$ nad zbiorem $F$ nazywamy funkcję $g : F \times X \rightarrow X$.

\subsection{Przykłady}
\begin{itemize}
    \item Mnożenie wektora przez skalar.
\end{itemize}

%%%%%%%%%%%%%%%%%%%%%%%%%%%%%%%%%%%%%%%%%%%%%%%%%%%%%%%%%%%%%%%
\section{Grupa}
Para $(X, *)$, gdzie $X \neq \emptyset$ a $*$ działaniem na tym zbiorze, jest grupą jeśli spełnione są postulaty:
\begin{enumerate}
    \item $\forall_{x,y \in X} x * y \in X$ - działanie wewnętrzne,
    \item $\forall_{x, y, z \in X} (x * y) * z = x * (y * z)$ - działanie łączne,
    \item $\exists_{e \in X} \forall_{x \in X} x * e = e * x = x$ - działanie posiada element neutralny,
    \item $\forall_{x \in X} \exists_{\tilde{x} \in X} x * \tilde{x} = \tilde{x} * x = e$ - każdy element posiada element odwrotny.
\end{enumerate}
Jeżeli dodatkowo spełniony jest postulat:
\begin{enumerate}
    \setcounter{enumi}{4}
    \item $\forall_{x,y \in X} x * y = y * x$ - działanie przemienne,
\end{enumerate}
to grupę nazywamy grupą przemienną (abelową).

W grupie $(X, *)$ potrafimy rozwiązywać równania (szukamy $x$):
\begin{itemize}
    \item $a * x = b \implies x = \tilde{a} * b$,
    \item $x * a = b \implies x = b * \tilde{a}$.
\end{itemize}

\subsection{Intuicja}
Grupa to zestaw elementów, które potrafimy dodawać i odejmować.

\subsection{Przykłady}
\begin{itemize}
    \item $(\{-1, 1\}, \cdot)$ - multiplikatywna grupa abelowa
    \item $(\mathbb{R}, +)$ - addytywna grupa abelowa
    \item $(\mathbb{R}, \cdot)$ - nie jest grupą, ponieważ 0 nie ma elementu odwrotnego
    \item $(\mathbb{R} - \{0\}, \cdot)$ - multiplikatywna grupa abelowa
    \item $(\{0, 1, 2, 3\}, \bigoplus)$, gdzie $a \bigoplus b = mod_{4}(a + b)$ - addytywna grupa abelowa
\end{itemize}

\subsection{Przykład (prawie)praktyczny - kodowanie danych}
Weźmy alfabet 26-literowy i ponumerujmy litery:
\[
    A=0,\ B=1,\ldots,Z=25.
\]
Zbiór:
\[
    \mathbb Z_{26}=\{0,1,\ldots,25\}
\]
z dodawaniem modulo 26 tworzy grupę:
\[
    (\mathbb Z_{26},+_{26}).
\]

\subsubsection{Szyfrowanie}
Ustalmy klucz szyfrujący:
\[
    k=3.
\]
Szyfrowanie polega na:
\[
    E_k(x)=x+_{26}k.
\]
Czyli:
\[
    A\rightarrow D,\quad
    B\rightarrow E,\quad
    C\rightarrow F,\ldots
\]
Przykładowo: $KOT$ zamieniamy na
\[
    10,\ 14,\ 19
\]
i dodajemy 3:
\[
    13,\ 17,\ 22
\]
czyli: $NRW$

\subsubsection{Odszyfrowanie}
Potrzebujemy elementu odwrotnego $-k$:
\[
    D_k(y)=y+_{26}(-k).
\]
W języku teorii grup:
\[
    -k=23\pmod{26}.
\]
Czyli klucz szyfrujący to element grupy
\[
    3\in\mathbb Z_{26},
\]
a jego element odwrotny to
\[
    3^{-1}_{\text{grupowo}}=-3=23\pmod{26}.
\]
Nie chodzi tutaj o odwrotność multiplikatywną, tylko o element odwrotny w~grupie
$(\mathbb Z_{26},+_{26})$. 
Mamy więc:
\[
    E_3(x)=x+_{26}3
\]
oraz
\[
    E_{-3}(x)=x+_{26}(-3).
\]
A następnie:
\[
    E_{-3}(E_3(x)) = (x+_{26}3)+_{26}(-3) = x.
\]
To dokładnie jest wykorzystanie elementu odwrotnego w grupie.

Co się stanie, jeśli najpierw przesuniemy litery o 3, a potem o 7?
Otrzymujemy:
\[
    (x+_{26}3)+_{26}7=x+_{26}10.
\]
Czyli dwa szyfrowania możemy połączyć w jedno:
\[
    E_7\circ E_3=E_{10}.
\]
A to jest właśnie działanie grupowe:
\[
    \boxed{
    \text{klucz }3
    \quad+\quad
    \text{klucz }7
    \quad=\quad
    \text{klucz }10
}
\]
wszystko modulo 26.

%%%%%%%%%%%%%%%%%%%%%%%%%%%%%%%%%%%%%%%%%%%%%%%%%%%%%%%%%%%%%%%
\section{Ciało}
Trójka $(X, \bigoplus, \bigotimes)$, gdzie $X \neq \emptyset$, a $\bigoplus$ i $\bigotimes$ są działaniami wewnętrznymi na zbiorze $X$, nazywana jest ciałem jeżeli spełnione są postulaty:
\begin{enumerate}
    \item $\forall_{x,y \in X} x \bigoplus y \in X$,
    \item $\forall_{x, y, z \in X} (x \bigoplus y) \bigoplus z = x \bigoplus (y \bigoplus z)$,
    \item $\exists_{\textbf{0} \in X} \forall_{x \in X} x \bigoplus \textbf{0} = \textbf{0} \bigoplus x = x$,
    \item $\forall_{x \in X} \exists_{-x \in X} x \bigoplus (-x) = (-x) \bigoplus x = \textbf{0}$,
    \item $\forall_{x,y \in X} x \bigoplus y = y \bigoplus x$,
\end{enumerate}
(postulaty 1-5 można zapisać krótko jako: $(X, \bigoplus)$ jest grupą abelową)
\begin{enumerate}
    \setcounter{enumi}{5}
    \item $\forall_{x,y \in X} x \bigotimes y \in X$,
    \item $\forall_{x, y, z \in X} (x \bigotimes y) \bigotimes z = x \bigotimes (y \bigotimes z)$,
    \item $\exists_{\textbf{1} \in X} \forall_{x \in X} x \bigotimes \textbf{1} = \textbf{1} \bigotimes x = x$,
    \item $\forall_{x \in X-\{\textbf{0}\}} \exists_{x^{-1} \in X} x \bigotimes x^{-1} = x^{-1} \bigotimes x = \textbf{1}$,
\end{enumerate}
(postulaty 6-9 można zapisać krótko jako: $(X - \{\textbf{0}\}, \bigotimes)$ jest grupą)
\begin{enumerate}
    \setcounter{enumi}{9}
    \item $\forall_{x,y,z \in X} (x \bigoplus y) \bigotimes z = (x \bigotimes z) \bigoplus (y \bigotimes z)$ oraz $z \bigotimes (x \bigoplus y) = (z \bigotimes x) \bigoplus (z \bigotimes y)$ - rozdzielność $\bigotimes$ względem $\bigoplus$.
\end{enumerate}
Jeżeli ponadto zachodzi:
\begin{enumerate}
    \setcounter{enumi}{10}
    \item $\forall_{x,y \in X} x \bigotimes y = y *\bigotimes x$.
\end{enumerate}
to ciało nazywamy przemiennym.

Intuicyjnie można powiedzieć, że ciało jest zbiorem, na którym potrafimy "dodawać" i "mnożyć".

\subsection{Przykłady}
\begin{itemize}
    \item $(\mathbb{W}, +, \cdot)$ - przemienne ciało liczb wymiernych,
    \item $(\mathbb{R}, +, \cdot)$ - przemienne ciało liczb rzeczywistych.
\end{itemize}

\subsection{Przykład (prawie)praktyczny - interpolacja wielomianowa w \(\mathbb Z_7\)}
Mamy trzy punkty:
$$ (1,2),\qquad(2,4),\qquad(4,1) $$
i szukamy wielomianu stopnia co najwyżej 2:
$$ f(x)=ax^2+bx+c $$
takiego, że:
$$ f(1)=2,\qquad f(2)=4,\qquad f(4)=1 $$
Wszystkie rachunki wykonujemy modulo 7.

\begin{enumerate}
    \item Układ równań
    Podstawiamy punkty:
    $$ a+b+c=2 $$ $$ 4a+2b+c=4 $$ $$ 16a+4b+c=1. $$
    Redukujemy modulo 7:
    $$ \begin{cases} a+b+c=2\\ 4a+2b+c=4\\ 2a+4b+c=1 \end{cases} $$
    Odejmujemy pierwsze równanie od drugiego:
    $$ 3a+b=2. $$
    Od trzeciego odejmujemy pierwsze:
    $$ a+3b=-1\equiv6. $$
    Mamy więc:
    $$ \begin{cases} 3a+b=2\\ a+3b=6 \end{cases} $$

    \item Wykorzystujemy odwrotność w ciele
    Z pierwszego:
    $$ b=2-3a. $$
    Podstawiamy:
    $$ a+3(2-3a)=6. $$ $$ a+6-9a=6. $$
    Modulo 7:
    $$ -8a=0. $$
    Ponieważ:
    $$ -8\equiv6\pmod7, $$
    mamy:
    $$ 6a=0, $$
    czyli:
    $$ a=0. $$
    Stąd:
    $$ b=2. $$
    A z pierwszego równania:
    $$ c=0. $$
    Otrzymujemy:
    $$ \boxed{f(x)=2x\pmod7}. $$
    Sprawdzenie:
    $$ f(1)=2, $$ $$ f(2)=4, $$ $$ f(4)=8\equiv1\pmod7. $$
    Działa.
\end{enumerate}

\end{document}